\paragraph{F.49 (Proved). Further field ranges where pinned reconstruction fails.}
The functor from finite-type $k$-schemes with universal homeomorphisms
inverted to pinned etale topoi over $B G_k$ is not full in each of
the following cases:
\begin{enumerate}
\item $k$ is any algebraic extension of a finite field, including
      $\overline{\mathbb F}_p$;
\item $k$ is infinite and perfect and $Z(G_k)\ne1$; in particular,
      $k$ is any real closed field;
\item $k$ is any algebraically closed field.
\end{enumerate}
In every case an endomorphism of $(\mathbb A^1_k)_{\mathrm{et}}$
already gives a counterexample. The relative comparison with the
base topos is part of the morphism. Forgetting that comparison would
erase the counterexamples constructed in the second case.

\emph{Manuscript problem locus.} Montpellier cote149, \S2,
fol.25v, scan52, including the margin asking about characteristic
zero; these statements sharpen the field-range delimitation in F.44.

\emph{Proof: the central comparison twist and its point function.}
Let $k$ be perfect, choose a separable closure $K$, put $H=G_k$,
$X=\mathbb A^1_k$, and let $p:X_{\mathrm{et}}\to BH$ be the
structure morphism. A central element $z\in H$ acts naturally and
$H$-equivariantly on every discrete continuous $H$-set. Pulling back
this natural automorphism to $X_{\mathrm{et}}$ gives an invertible
two-cell of $p$. Define
\[
 T_z=(\operatorname{id}_{X_{\mathrm{et}}},\alpha_z)
       \quad\text{in the category of topoi over }BH,
 \qquad \alpha_z:p\circ\operatorname{id}\Longrightarrow p.
                                                               \tag{F49.1}
\]
Choose the orientation of this comparison so that on the fibre of a
constant finite etale scheme $E/k$ it is
\[
 \operatorname{Hom}_k(E,K)\longrightarrow
 \operatorname{Hom}_k(E,K),\qquad u\longmapsto z\circ u.
                                                               \tag{F49.2}
\]
The underlying geometric morphism is the identity, hence is pinned.

Its induced function on framed geometric points is
\[
                         a\longmapsto z(a),\qquad a\in K.
                                                               \tag{F49.3}
\]
Here is an explicit check of the relative data. A standard geometric
point at $a$ has the standard identification of its restriction to
constant etale schemes with the fixed fibre functor of $BH$.
Composition with (F49.1) changes that identification to (F49.2).
Applying $z$ to every lift of the point through an etale neighbourhood
sends a lift above $a$ to a lift above $z(a)$. On constant covers
this is exactly (F49.2), so it identifies the twisted framed point
with the standard framed point at $z(a)$.

Distinct elements of $K$ give nonisomorphic relative closed geometric
points. For completeness, over a closed point with residue field
$E\subset K$ their groupoid is the homotopy fibre of
$B G_E\to B G_k$ at the fixed fibre functor. Since $G_E\to G_k$
is injective, this groupoid has trivial automorphism groups, and its
isomorphism classes are the embeddings $E\hookrightarrow K$.
Different closed supports have no isomorphisms between them.
This is also the case $L=K$ of Carlson--Haine--Wolf,
Proposition~3.6 \cite{CHW2024}. Consequently (F49.3) distinguishes
the relative morphisms $T_z$ without any assertion about the centre
of the generic Galois group of $X$.

\emph{Proof: infinite perfect fields with nontrivial Galois centre.}
Take $1\ne z\in Z(G_k)$. Formula (F49.3) fixes every $a\in k$
and moves some $a\in K$. No endomorphism from the scheme
localization can have this point function. We make this assertion
explicit, including the inverted universal homeomorphisms.
By the absolute-weak-normalization description of that localization
(F.4; Carlson--Haine--Wolf, Proposition~2.28), its endomorphisms
from $\mathbb A^1_k$ to itself are represented by
\[
 \operatorname{Hom}_k((\mathbb A^1_k)^{\mathrm{awn}},\mathbb A^1_k).
                                                               \tag{F49.4}
\]
In characteristic zero $\mathbb A^1_k$ is already absolutely weakly
normal, so a member is given by $P(t)\in k[t]$. If $P(a)=a$ for
every $a\in k$, the infinitude of $k$ gives $P=t$.
In characteristic $p$, (F49.4) is the set
\[
                  k[t^{1/p^\infty}]
                    =\bigcup_{n\geq0}k[t^{1/p^n}].
\]
For an element $h$ of this ring choose $n$ so that
$h^{p^n}=P(t)\in k[t]$. If $h(a)=a$ for every $a\in k$, then
$P(a)=a^{p^n}$ on an infinite set; hence $P=t^{p^n}$.
The ring is reduced, so Frobenius is injective and $h=t$.
In either characteristic the only possible scheme morphism is the
identity. Its geometric-point function fixes all of $K$, whereas
(F49.3) does not. Thus $T_z$ is an explicit morphism outside the
image. A real closed field has algebraic closure $k(i)$ and
$G_k=C_2$; complex conjugation supplies the stated example.

\emph{Proof: all algebraic extensions of finite fields.}
Let $k\subset\overline{\mathbb F}_p$ be algebraic over
$\mathbb F_p$ and write
\[
 H=\operatorname{Gal}(\overline{\mathbb F}_p/k)
       \subset\operatorname{Gal}(\overline{\mathbb F}_p/\mathbb F_p)
       =\widehat{\mathbb Z}.
\]
If $H\ne1$, it is abelian and has cardinality $2^{\aleph_0}$.
Indeed a nonzero element has a nonzero projection to some
$\mathbb Z_\ell$; the closure of its integer multiples projects
onto a nonzero subgroup of $\mathbb Z_\ell$, already of that
cardinality. Every $z\in H$ gives (F49.1), and (F49.3) makes
their isomorphism classes distinct. On the scheme side there are
only countably many localized endomorphisms: $k$ is countable and
the ring $k[t^{1/p^\infty}]$ in (F49.4) is countable. Thus the
functor cannot be full. This handles finite fields as well as all
proper infinite algebraic subfields. For the infinite subfields the
preceding argument also gives the stronger explicit obstruction for
every nonidentity $z\in H$.

If $H=1$, then $k=\overline{\mathbb F}_p$ and $BH=\mathrm{Set}$.
For each $\sigma\in\widehat{\mathbb Z}$ the ring automorphism
\[
 k[t]\longrightarrow k[t],\qquad
    \sum_i a_i t^i\longmapsto\sum_i\sigma^{-1}(a_i)t^i
                                                               \tag{F49.5}
\]
gives a semilinear automorphism of the underlying scheme.
Its map on closed points is $a\mapsto\sigma(a)$.
The resulting topos automorphism is automatically relative to
$BH=\mathrm{Set}$, and it is pinned. Distinct $\sigma$ induce
distinct point maps, so they give $2^{\aleph_0}$ distinct morphisms
again, compared with the same countable ring (F49.4). This proves
the first assertion, including the algebraically closed endpoint.

\emph{Proof: arbitrary algebraically closed fields.}
Let $k$ now be any algebraically closed field. It suffices to find
a nonidentity field automorphism $\sigma$ fixing an infinite subfield
$k_0\subset k$. The semilinear construction (F49.5), relative to
$\mathrm{Set}$, has point function $a\mapsto\sigma(a)$. It fixes
all $k_0$-points and moves another point. Exactly the polynomial or
finite-$p$-denominator calculation used above shows that a morphism
in (F49.4) fixing this infinite subset must be the identity.
Therefore (F49.5) is outside the image.

In characteristic zero, choose a copy of $\overline{\mathbb Q}$ in
$k$, take a nonidentity automorphism of it fixing $\mathbb Q$, and
extend it to $k$. Such extension is obtained by fixing a
transcendence basis over $\overline{\mathbb Q}$ and extending the
resulting automorphism of the rational-function field to its
algebraic closure. The fixed subfield contains $\mathbb Q$.
In characteristic $p$, choose
$\delta\in\widehat{\mathbb Z}=\prod_\ell\mathbb Z_\ell$ whose
$2$-adic coordinate is zero and whose odd-primary coordinates are
one. The automorphism $\operatorname{Frob}_p^\delta$ of
$\overline{\mathbb F}_p$ is nonidentity and fixes the infinite field
\[
                    \bigcup_{n\geq0}\mathbb F_{p^{2^n}}.
\]
Extend it to $k$ by the same transcendence-basis argument. This
supplies $\sigma$ and $k_0$ in every characteristic, and completes
all three assertions. The attached program \texttt{code/F49\_checks.py}
checks the framed-point distinction over $\mathbb F_9$ and 2726 finite
compatibilities of the displayed profinite exponent; the infinitary
claims follow from the proofs above. \hfill$\square$
